% =====================================================================
% RASCUNHOS para os pontos do orientador (REESCREVER na sua voz).
% 1) Apresentação do app (Seção 5.1) com screenshots
% 2) Seção "Dificuldades encontradas"
% 3) Parágrafo de limitação sobre a fluidez do PWA
% ATENÇÃO: são rascunhos — reescreva com suas palavras (política de IA +
% checagem de plágio). Confira números e insira as capturas de tela reais.
% =====================================================================

% ---------------------------------------------------------------------
% (1) SEÇÃO 5.1 — APRESENTAÇÃO DO APLICATIVO (substitui/expande a atual)
% Tire 4 capturas de tela no app/PWA e salve no projeto do Overleaf.
% ---------------------------------------------------------------------
\subsection{O produto: o aplicativo InterpreteLabBR}
\label{sub:produto}

O artefato resultante é o cliente móvel do InterpreteLabBR, desenvolvido em React
Native/Expo e consumindo a mesma API FastAPI do PWA. O fluxo de uso é direto. Na
tela inicial, o usuário informa sexo e idade e escolhe entre dois modos de entrada:
a digitação manual dos valores do hemograma (Figura~\ref{fig:tela-manual}) ou o
envio do laudo em PDF, do qual o sistema extrai automaticamente os analitos
(Figura~\ref{fig:tela-pdf}). Após a análise, o aplicativo exibe a tela de
resultados, com a classificação de cada analito segundo a referência da PNS, as
especialidades sugeridas e um \emph{briefing} educativo ao paciente
(Figura~\ref{fig:tela-resultados}). O sistema apresenta ainda uma tabela que
compara, lado a lado, a classificação pela referência da PNS e pela referência
clássica impressa no laudo, sinalizando divergências
(Figura~\ref{fig:tela-comparacao}). Essa apresentação contextualiza as telas e as
interações que são discutidas nas seções seguintes.

% Substitua os 4 arquivos de imagem pelos nomes das suas capturas de tela.
\begin{figure}[H]\centering
\includegraphics[width=0.32\textwidth]{tela_manual.png}
\caption{Entrada manual dos valores do hemograma.}\label{fig:tela-manual}
\end{figure}
\begin{figure}[H]\centering
\includegraphics[width=0.32\textwidth]{tela_pdf.png}
\caption{Envio do laudo em PDF para extração automática.}\label{fig:tela-pdf}
\end{figure}
\begin{figure}[H]\centering
\includegraphics[width=0.32\textwidth]{tela_resultados.png}
\caption{Tela de resultados, com classificação, especialidades e \emph{briefing}.}\label{fig:tela-resultados}
\end{figure}
\begin{figure}[H]\centering
\includegraphics[width=0.32\textwidth]{tela_comparacao.png}
\caption{Comparação entre a referência da PNS e a referência clássica do laudo.}\label{fig:tela-comparacao}
\end{figure}

% ---------------------------------------------------------------------
% (2) NOVA SEÇÃO — DIFICULDADES ENCONTRADAS (coloque perto da Conclusão)
% ---------------------------------------------------------------------
\section{Dificuldades Encontradas}
\label{sec:dificuldades}

O desenvolvimento e a avaliação do artefato enfrentaram percalços que exigiram
ajustes de rota. Na extração de dados, o \emph{parser}, originalmente validado para
o leiaute do SUS, falhava ao processar a série vermelha de laudos de rede privada,
cujo texto extraído apresentava o nome do analito colado ao valor; o problema foi
diagnosticado por inspeção do texto extraído e corrigido tornando os padrões de
reconhecimento tolerantes à ausência de espaço. Na Frente~A, a medição da fluidez
de rolagem revelou uma assimetria entre os clientes: o \texttt{dumpsys gfxinfo} não
captura os quadros do Chrome (que utiliza o \emph{pipeline} Blink), o que impediu
uma comparação direta de FPS; optou-se por reportar com rigor a métrica
comparável (memória) e documentar a limitação. Na Frente~C, a dificuldade de
reunir, no prazo, um painel de cinco especialistas levou à condução da inspeção
como estudo-piloto, com um painel reduzido. Houve ainda desafios operacionais, como
a geração e a distribuição do aplicativo aos avaliadores (via \emph{build} em nuvem
e instalação fora da loja oficial) e a sincronização entre a correção do
\emph{backend} e a sua publicação em produção. Esses contratempos, além de
superados, orientaram decisões metodológicas e ampliaram o aprendizado ao longo da
pesquisa.

% ---------------------------------------------------------------------
% (3) LIMITAÇÃO — FLUIDEZ DO PWA (acrescente na Seção 5.7 / limitações)
% ---------------------------------------------------------------------
% Parágrafo para colar junto às limitações da Frente A:
Reconhece-se como limitação relevante a ausência de medição da fluidez da interface
do PWA. O instrumento empregado (\texttt{dumpsys gfxinfo}) não expõe os quadros do
Chrome; existiriam, contudo, alternativas para aferi-la --- por exemplo, as
\emph{DevTools} do Chrome, cujas métricas de desempenho podem ser coletadas
remotamente via ADB; a API \texttt{requestAnimationFrame} para telemetria em
JavaScript; e ferramentas dedicadas como o Perfetto, que cobrem tanto o PWA quanto
o cliente nativo. A não adoção dessas alternativas decorreu do escopo e do prazo
desta etapa, concentrados na comparação já viável (memória e inicialização), e
constitui direção imediata de trabalho futuro.
